% compendium/international-equity-markets/equity-indices.tex
% Entry migrated on 2026-10-16 from the legacy macro chapter (slice 37, 02_transmission_maps.tex), owned by the
% section overview international-equity-markets. Gliner (2014) is cited by chapter, section and PDF page of the ebook.
% Dispositions in notes/migration-slices.md, slice 37.

\section{Equity indices and instruments for macro}

\begin{qfdefinition}[Equity indices and instruments for macro: major indices by region, price, capitalisation and equal weighting, and ETFs, ADRs, index futures and options as ways to trade them]\label{def:equity-indices-derivatives}
\qffield{Statement} Gliner's Table~8.1 lists major indices by region: United States (S\&P 500, Nasdaq, DJIA, Russell 2000, Wilshire 5000), Europe (Euro Stoxx 50, DAX, CAC 40, FTSE 100, FTSE MIB, IBEX 35, AEX), Asia ex-Japan (Shanghai Composite, Hang Seng, KOSPI, ASX 200, NZX 50), Japan (Nikkei, Topix, JASDAQ) and emerging markets (Bovespa, NIFTY, Micex and RTS, Mexico IPC, TOP 40 and JALSH, BIST 30 and 100, Jakarta Composite); indices ``can be traded directionally or as a relative value trade'', and their exposures differ, so ``the Bovespa is more likely to be affected by commodities than the NIFTY, while the S\&P 500 has more financial exposure than the Nasdaq, which is more technology-weighted'' \citep[chap.~8, sec.\ ``Equity Indices Overview'', PDF pp.~214--215]{s7-gliner}. Price-weighted indices consider only prices and ``can be strongly affected by the move of one'' stock; capitalisation-weighted indices are dominated by the largest companies; ``it is useful to know the weighting structure of each index in order to accurately interpret its meaning'' (pp.~214--215). Instruments (sec.\ ``Equity Derivatives'', PDF pp.~218--221): ETFs, ``a fast, cheap way to get exposures'' that trade like stocks but charge fees, whose composition can differ from the index (EWZ held as much as 15 percent Petrobras, more than the Bovespa); ADRs, a depositary bank's holdings of a non-U.S.\ stock listed in the United States in dollars, ``subject to foreign exchange risk''.

\qffield{Notation} Price-weighted: $I_t\propto\sum_i P_{i,t}$; capitalisation-weighted: $I_t\propto\sum_i N_{i,t}P_{i,t}$ (often free-float adjusted); equal-weighted: equal weights reset at each rebalance.

\qffield{Conditions} Index lists, constituents and ETF holdings are as of 2013. \emph{Two changes checked} (second slice-37 follow-up). JASDAQ is no longer a market segment: ``On April 4, 2022, TSE reorganized the former market segments (First Section, Second Section, Mothers and JASDAQ) into three new market segments: Prime, Standard and Growth'' \citep[market-structure page, retrieved 2026-10-16]{jpx-market-restructure}, so a JASDAQ index in Gliner's sense describes a market that no longer exists. Moscow Exchange now publishes the ``MOEX Russia Index and RTS Index'', listing ``the MICEX Index'' only as a registered trademark \citep[index page, retrieved 2026-10-16]{moex-indices}; Gliner's ``Micex'' is the MOEX Russia Index (the date of the renaming was not given on the page). \emph{Two weighting schemes checked} (source follow-up, 2026-10-16). The Nikkei 225 ``is a price-weighted equity index, which consists of 225 stocks in the Prime Market of the Tokyo Stock Exchange'', reviewed semi-annually and computed every five seconds \citep[profile page, retrieved 2026-10-16]{nikkei-225-profile}; TOPIX ``is a free-float adjusted market capitalization-weighted index'', with base value 100 on 4 January 1968 \citep[TOPIX page, retrieved 2026-10-16]{jpx-topix}. So Japan's two headline indices sit on opposite sides of the price-versus-capitalisation distinction above. \emph{Two more checked} (second source follow-up, 2026-10-16). South Africa: Gliner's TOP 40 and JALSH are the FTSE/JSE Top 40 and All Share indices, two of seven ``headline'' indices in FTSE Russell's FTSE/JSE Africa series; the Top 40 is ``the 40 most investable companies ranked by investable market capitalisation, in the FTSE/JSE All Share Index'', and the headline indices ``were harmonised effective 18 March 2024'' to reflect the SWIX free float, the shareholder-weighted method of the rules' Appendix~C \citep[v9.3, September 2026, secs.~1.7.1 and~4.5.1, pp.~4 and~13]{ftse-jse-ground-rules-2026}. United States: the Wilshire 5000 is now the FT Wilshire 5000, ``a comprehensive, float adjusted measure of the U.S. stock market, designed to reflect the performance of all publicly traded U.S. equity securities that have readily available prices'', with capitalisation-based size segments \citep[WI-METH-001, January 2026, sec.~1.1, p.~4]{ft-wilshire-5000-2026}; its name is not a constituent count. \emph{Two more} (batch of 2026-10-17). India: ``The Nifty 50 is a 50 stock, float-adjusted market-capitalization weighted index for India'', its stocks representing ``about 66\%'' of NSE's float-adjusted capitalisation \citep[Methodology Document for Equity Indices, September 2026, PDF p.~19]{nse-nifty-methodology-2026} (the same page's turnover sentence also says ``float-adjusted market capitalization'', apparently a slip). Turkey: BIST indices weight by free-float market value, with a weighting factor below one only in capped versions; BIST~100 ``automatically covers BIST 30 and BIST 50 stocks'', and the capped variants bound weights at 10 or 25 percent with thresholds of 15 or 30 percent \citep[Ground Rules, April 2020, secs.~2.2--2.3, 2.20 and~3.1, pp.~1--5]{bist-ground-rules-2020}, the edition served at Borsa Istanbul's address on 2026-10-17. Hong Kong: the Hang Seng Index is ``freefloat-adjusted market capitalisation weighted'' with the number of constituents ``fixed at 100''. Individual weights are capped at 8 percent for non-foreign constituents and 4 percent for foreign ones, 10 percent in aggregate. It is reviewed quarterly, launched 24 November 1969 with a base of 100 on 31 July 1964 \citep[Hang Seng Index and Sub-indexes methodology, version~2.3 of July 2023, PDF pp.~3, 5 and~6--7]{hsi-methodology-2023}. \emph{Five more} (third batch of 2026-10-17). The DAX, EURO STOXX~50 and DJIA rules were read from Internet Archive captures, because the providers' own sites refused the script or could not be verified. The Russell and FTSE rules came directly from FTSE Russell.
\begin{itemize}
\item The DJIA ``is a 30-stock, price-weighted index'', drawn from S\&P~500 members outside the transportation and utilities groups (GICS 2030 and~55) \citep[May 2026, ``Index Objective and Highlights'' and ``Eligibility Criteria'', PDF pp.~4--5; capture of 27 September 2026]{spdji-dj-averages-2026}. It is the price-weighted case of the notation below.
\item The Russell 2000 is companies ranked 1{,}001 to 3{,}000 by total market capitalisation, within ``float-adjusted, market capitalization-weighted'' indexes. Reconstitution is now semi-annual, ``on the fourth Friday in June and the second Friday in December'', with IPOs added quarterly \citep[v7.2, August 2026, secs.~1.1.1, 4.2.3 and~6.10, PDF pp.~3, 9 and~20]{russell-us-construction-2026}.
\item The FTSE~100 ``will consist of the largest 100 UK companies by full market capitalisation (i.e.\ before the application of any investability weightings)''. Constituents are then weighted with free-float and investability adjustments and reviewed quarterly \citep[v17.3, September 2026, secs.~6.4 and~7.1.1, PDF pp.~14 and~16]{ftse-uk-ground-rules-2026}.
\item The EURO STOXX~50 holds 50 Eurozone supersector leaders, weighted by free-float market capitalisation and capped at 10~percent quarterly. It is selected with a 40--60 buffer and reviewed annually in September \citep[Index Methodology Guide, October 2026, sec.~9.2, PDF pp.~100--101; capture of 16 October 2026]{stoxx-index-guide-2026}.
\item The DAX ``includes the 40 largest companies in terms of market capitalisation on the Regulated Market of FWB'' that meet quality and profitability requirements \citep[Guide to the DAX Equity Indices, version 11.3.0 of 3 August 2022, PDF p.~21; capture of 2 September 2025]{dax-guide-2022}. A later version may exist; none was captured.
\end{itemize}
Still not checked: the S\&P~500, ASX~200, NZX~50 and Mexico IPC (S\&P DJI's own PDFs returned HTTP 403, and the Archive's capture of the U.S. methodology returned 503 twice); Nasdaq (indexes.nasdaqomx.com, HTTP 403); CAC~40 and AEX (Euronext, 403, and no capture at the address tried); FTSE~MIB (404 at the address tried); KOSPI, the Shanghai Composite, Bovespa and the Jakarta Composite (no new route). \emph{IBEX~35} (fourth batch of 2026-10-17, from the index committee's notice of 12 December 2024 through an Internet Archive capture; BME's site served a web page in place of the PDF): constituents carry free-float factors in bands (the notice moves Colonial from 60 to 40 percent and Indra from 80 to 100 percent), and ``The individual weighting of stocks will be adjusted to fulfil the 20\% limit'', with a 40 percent limit in the bank, energy and construction sub-indices; changes applied from 23 December 2024, and only the Spanish text is legally binding \citep[Notice 18/2024, pp.~2--3]{bme-ibex-notice-18-2024}. So the IBEX~35 is free-float capitalisation-weighted with a 20 percent cap; BME's Technical Regulations themselves were not read. The earlier note that the Hang Seng address returned 404 is superseded: the HSI methodology was read in the second batch of 17 October 2026.

\qffield{Practice} For relative-value index trades, size legs by beta and volatility, and hedge or deliberately keep the currency exposure; for the section's outcome, compare markets in local currency and in dollars, hedged and unhedged.

\qffield{Pitfalls} An ADR's dollar price follows the local share price converted at the exchange rate, so it carries currency risk for a dollar investor, not because ``shares are held locally'' but because the underlying is priced in local currency. The legacy's index futures (``liquid, leveraged exposure'') and options (``asymmetric payoffs'') are generic; futures conventions are in \hyperref[conv:cme-globex-trading]{CME Globex trading} and ETF mechanics in \hyperref[mech:etf-creation-redemption]{ETF creation and redemption}.

\qffield{See also} \hyperref[meth:equity-index-valuation]{equity-index valuation}; \hyperref[def:vix-variance-dividend-swaps]{VIX and variance swaps}; \hyperref[def:global-macro-strategy]{global macro}.

\qffield{Further reading} The section's readings on market integration and cross-listing.

\qffield{Index} equity indices; S\&P 500; Euro Stoxx 50; Nikkei; Topix; emerging market indices; MSCI; price-weighted index; capitalisation-weighted index; equal-weighted index; relative value index trade; ETF; ADR; American depositary receipt; index futures

\qffield{Status} Verification: partly checked 2026-10-16 (Gliner's text, chap.~8, PDF pp.~214--221, with Tables~8.1--8.2; second slice-37 follow-up, 2026-10-16: the JPX market-structure page and Moscow Exchange's index page read; source follow-up of 2026-10-16: Nikkei's Nikkei 225 profile and JPX's TOPIX page; second source follow-up of 2026-10-16: FTSE/JSE ground rules v9.3 and the FT Wilshire 5000 methodology of January 2026; batch of 2026-10-17: NSE Indices' methodology of September 2026 and Borsa Istanbul's ground rules of April 2020, text layers; second pass: Hang Seng Indexes' HSI methodology v2.3 of July 2023, text layer; third batch of 2026-10-17: Dow Jones Averages methodology (May 2026), Russell US construction v7.2, FTSE UK ground rules v17.3, STOXX guide (October 2026) and DAX guide v11.3.0, text layers; fourth batch: BME's IBEX committee notice 18/2024). Stability: draft. Currency: index lists and ETF examples as of 2013; fourteen indices' names or weighting as of 2026 (BIST rules dated 2020, DAX 2022, HSI 2023).
\end{qfdefinition}
